% Source: sources/returns/S6_04_MAPS_MANIFOLDS/sections/manifolds.tex:2364–2599
\paragraph{Theorem D.2.12 (Proved: seven further unbounded one-log families).}
Let $G$ be the marked Mordell--Weil generator of the indicated
rational elliptic surface. In the table, the superscript $\bullet$
marks the fibre supporting the degree-one linearization, and $d$
is the linear order at the sole additive fibre.
\[
 \begin{array}{c|l|rr|r}
 \mathrm{OS}&\text{ordered fibre configuration}&P/G&Q/G&d\\\hline
 45&I_8,II,I_1^\bullet,I_1&8&1&6\\
 47&I_7,I_2,I_1^\bullet,II&14&1&6\\
 47&I_7,I_1^\bullet,I_1,III&7&2&4\\
 55&I_5,I_4,II,I_1^\bullet&20&1&6\\
 56&I_5,I_2,IV,I_1^\bullet&10&3&3\\
 56&I_5,I_3,III,I_1^\bullet&15&2&4\\
 49&I_2,IV^*,I_1^\bullet,I_1&2&3&3
 \end{array}
\]
Use $M=\mathcal O(P-O)$ and translation by $Q$.
Keep the two weight-zero multiplicative fillings and the
primitive toroidal filling at $I_1^\bullet$.
In the integral marking specified below, put
\[
 m=dk,\qquad v=n\ell+r\delta,
 \qquad k>0,\quad\gcd(n,dk)=1,
 \qquad n,r\in\mathbb Z.
\]
Every row admits these logarithmic fillings for every displayed
integer choice, and the resulting compact smooth threefold has
\[
 \begin{gathered}
 \pi_1(X)=\mathbb Z/|n|,\qquad
 H_i(X;\mathbb Z)=\mathbb Z/|n|\quad(1\leq i\leq4),\\
 H_5(X)=0,\qquad H_0(X)=H_6(X)=\mathbb Z.
 \end{gathered}
\]
In particular $n=\pm1$ gives $S^6$ for every $k>0$ in all seven
rows. Within the one-log construction, any simply connected
choice of affine vector must have $n=\pm1$, so the coprimality
hypothesis discards none of its sphere cases.
\begin{proof}
\emph{Geometric data and the actual multiplicative fillings.}
The configurations and height lattices are the primary
Oguiso--Shioda/Fukae data verified in D.2.1. Each displayed
pair has height pairing one. At every fibre at least one of
$P,Q$ is locally narrow, and at each weight-zero $I_j$ the
section $P$ is narrow while $Q$ generates its component group.
In the marked component groups the latter residues are
\[
 \begin{array}{c|l}
 \text{row}&Q\text{ on its higher multiplicative fibres}\\\hline
 1&3\pmod8\\
 2&5\pmod7,\quad1\pmod2\\
 3&3\pmod7\\
 4&2\pmod5,\quad3\pmod4\\
 5&3\pmod5,\quad1\pmod2\\
 6&2\pmod5,\quad1\pmod3\\
 7&1\pmod2.
 \end{array}
\]
All residues are units. At the additive fibre $Q$ is narrow;
its component residue is zero. The exact $P$ residues as well
as these $Q$ residues are retained with the matrices in the
accompanying \texttt{manifolds\_rankone\_families.json} record.

The line $J=M\otimes t_Q^*M^{-1}$ has multidegree zero on
every fibre: either $t_Q$ preserves its components, or $M$
already has multidegree zero there. Translation acts trivially on the degree-zero Picard group of
the generic elliptic fibre, so $J$ is generically trivial. Its
vertical divisor and zero component degrees then make it a
pullback from the base; all original fibres have multiplicity one.
Its degree is $\langle P,Q\rangle=1$, as all local mixed
height corrections vanish. Choose its section with its simple
zero at the marked $I_1^\bullet$. The contracting-quotient
construction and its properness estimate are the same as in
Engel's construction.

At a weight-zero $I_j$, the restriction of $M$ to each normalized
component $\mathbb P^1$ has degree zero and is trivial.
Translation by $Q$ cyclically permutes all $j$ components,
because its component residue is a unit. The quotient therefore
has one normalization component and one double curve.
After $j$ steps the first component returns; its normalization
is the quotient of $\mathbb P^1\times\mathbb C^*$ by a
translation on the second factor and a constant multiplicative
automorphism on the first. It is a $\mathbb P^1$ bundle over an
elliptic curve defined by a degree-zero line, hence is
topologically trivial. Its zero and infinity sections are
identified by a translation and flat units. These are isotopic
to the identity in the marked circle directions.
Each node has the original semistable chart $xy=t$ up to a
holomorphic unit: the infinite contracting action has no point
stabilizer and does not ramify this chart.
Thus the filled neighborhood and its normal-circle chains are
the primitive one-node model $N_{\rm node}\times T^2$ used
in D.2.6. In particular all three fixed one-cycles supply its
three normal tori, not just the vanishing circle.
This argument identifies the geometric local model, rather
than inferring it from a rational monodromy matrix.

At the marked $I_1$, the simple zero gives the usual one-component
hexagonal toroidal filling. Its two angular normal tori are
primitive. At the additive fibre, take the good-reduction root
model and then its degree-$k$ base change. An invariant vector
with primitive invariant-coordinate value $n$ defines a free
order-$dk$ affine action because $\gcd(n,dk)=1$.
The marked logarithmic collar and compact gluing are exactly
those proved in D.2.6. These steps prove existence and properness
for every allowed $k,n,r$.

\emph{Exact integral certificates, with no parameter bounds.}
Use the geometric open-family monodromy lifts from the section
cocycles and degree-one linearization, as in D.2.1 and the
public primary monodromy derivation of \cite{manExplorer26}.
Cyclically reorder the four matrices so that the additive
one is first, and denote them by $A_0,A_1,A_2,A_3$, with
$A_0A_1A_2A_3=I$.
The marking here is $(e_1,e_2,\delta,\gamma)$, with
$\chi=\gamma^*$. The following integral vectors and constants
are obtained by exact unimodular reduction:
\[
 \begin{array}{c|r|r|l}
 \text{row}&d&q&\ell\\\hline
 1&6&1&(-1,0,0,1)\\
 2&6&1&(0,-1,0,1)\\
 3&4&2&(-3,-2,0,1)\\
 4&6&1&(0,0,0,1)\\
 5&3&3&(-1,3,0,1)\\
 6&4&2&(-1,3,0,1)\\
 7&3&3&(-3,-1,0,1).
 \end{array}
\]
In every row $A_0\ell=\ell$, $A_i\delta=\delta$ and
$\chi(\ell)=1$.
All four matrices, integral fixed bases and the integer
left inverses used next are written in the retained JSON;
the executable derives and verifies them using only the
MIT monodromy source, without a local-model database or API.

Let $D=[A_0-I\ A_1-I\ A_2-I]$ and set
\[
 j_0(x)=(x,0,0),\quad j_1(x)=(0,x,0),\quad
 j_2(x)=(0,0,x),\quad
 j_3(x)=-(A_1A_2x,A_2x,x).
\]
For each $i=1,2,3$, choose the full integral fixed basis
$F_i=\ker(A_i-I)$, of rank two at the cusp and rank three
at a one-node fibre. Form the $12\times9$ matrix
\[
 K=[j_1(F_1)\ j_2(F_2)\ j_3(F_3)\ j_0(\ell)].
\]
The exact certificates give
\[
 DK=0,\quad\operatorname{rank}D=3,\quad LK=I_9,\quad
 (Lj_0\delta)_9=0,
 \tag{D.R1}
\]
where $L$ is the explicitly retained integer left inverse.
Thus the nine actual sweep tori form a primitive integral
basis of the graph kernel. The exterior coinvariant matrix
$[\bigwedge^2A_i-I]_{i=0}^3$ has five nonzero Smith entries,
all one. Its primitive dual row is the invariant form
\[
 \omega=P_0\,e_1^*e_2^*-Q_0\,\delta^*\gamma^*,
 \qquad P=P_0G,\quad Q=Q_0G.
 \tag{D.R2}
\]
The regular-family $H_2$ is therefore free of rank ten:
one fibre coinvariant of $\omega$-value one and the nine
sweeps. The constant form $\omega$ has zero value on the
latter actual fixed-loop sweeps, so this is an integral
splitting. Quotienting by the eight normal tori of the
unchanged fillings gives
$\mathbb Z\langle s,t\rangle$, with $t=j_0(\ell)$.
Equation (D.R1) also proves $j_0(\delta)=0$ in this quotient.

In the vertical lattice $\ker\chi$, the finite coinvariants
are free of rank one. Let $b$ be a primitive representative;
the exact fixed-basis computation gives
\[
 [\delta]=q[b],\qquad
 \mathcal N_d b=(d/q)\delta,\qquad
 \mathcal N_d\ell=d\ell,\qquad
 -\omega(b,\ell)=1,\quad
 -\omega(\ell,\delta)=-q.
 \tag{D.R3}
\]
Here $\mathcal N_d=\sum_{i=0}^{d-1}A_0^i$; all are
identities of integral vectors, not rational averages.
The finite reduced fibre has free $H_1$, with generators
$b,z$, where
\[
 mg=n\ell+rq b,\qquad
 z=\alpha\ell+\beta g,\qquad m\alpha+n\beta=1.
\]
Apply the normal-circle torus formula of D.2.6 to these two
cycles. The first column is the fibre area
$-\omega(b,v)$ and the sweep $j_0(\mathcal N_m b)$;
the second is $-\alpha\omega(\ell,v)$ and
$j_0(\alpha\mathcal N_m\ell+\beta v)$.
After the eight unchanged normal relations, (D.R1)--(D.R3)
give exactly
\[
 \begin{pmatrix}n&-qr\alpha\\0&1\end{pmatrix}.
 \tag{D.R4}
\]
Its Smith form is $(1,|n|)$ for all the asserted integers.

For van Kampen, the images of the three unchanged vanishing
lattices already generate the full primitive $\ker\chi$;
their combined matrix has Smith factors $(1,1,1)$, verified
in the same exact computation. Their meridians and the base
relation kill the remaining finite meridian. The only remaining
relation is $n\gamma=0$. This proves the nonabelian
fundamental-group assertion. It also shows, without assuming
coprimality, that any simply connected affine choice must have
$n=\pm1$.

The degree-one Mayer--Vietoris map has free source and target
of rank fourteen and finite cokernel $\mathbb Z/|n|$; hence
its kernel is zero. Every local degree-two boundary map is
onto, by the product models, the finite free quotient, and
the toroidal specialization proved above. Therefore (D.R4)
is the full remaining degree-two presentation and
$H_2(X)=\mathbb Z/|n|$.
Only the hexagonal cusp contributes Euler characteristic,
giving $\chi(X)=2$. Duality supplies all the other groups,
and homotopy-sphere recognition with $\Theta_6=0$ proves
the sphere statement.

The exact program \texttt{code/manifolds\_rankone\_families.py}
checks all constant matrices, integral bases and (D.R1)--(D.R3),
and also 234 full $10\times10$ Smith forms at displayed
sample parameters. Formula (D.R4) proves the result for
unrestricted orders and twists; the samples are corroboration.
No priority claim is made.
\end{proof}
